\documentclass[11pt]{amsart}
\usepackage[T1]{fontenc}
\usepackage[margin=1in]{geometry}
\usepackage{amsmath,amssymb,amsthm,booktabs,array}
\usepackage[colorlinks=true,linkcolor=blue,citecolor=blue,urlcolor=blue]{hyperref}
\newtheorem{theorem}{Theorem}[section]
\newtheorem{lemma}[theorem]{Lemma}
\newtheorem{corollary}[theorem]{Corollary}
\newtheorem{proposition}[theorem]{Proposition}
\theoremstyle{remark}
\newtheorem{remark}[theorem]{Remark}
\newcommand{\C}{\mathcal C}
\newcommand{\D}{\mathcal D}
\newcommand{\K}{\mathcal K}
\newcommand{\E}{\mathcal E}
\newcommand{\Q}{\mathcal Q}
\newcommand{\FPdim}{\operatorname{FPdim}}
\newcommand{\Hom}{\operatorname{Hom}}
\newcommand{\SU}{\mathrm{SU}}
\newcommand{\wt}[1]{\mathsf{#1}}
\newcommand{\one}{\mathbf 1}
\title[A rank-ten conformal-embedding model]{A unitary categorification of a rank-ten multiplicity-free fusion ring}
\date{24 October 2026}
\begin{document}
\begin{abstract}
We identify the census fusion ring \texttt{R10\_1\_189} with the
Grothendieck ring of a fusion subcategory arising from the known conformal
embedding $\SU(4)_6\subset\SU(10)_1$. A short proof uses two exterior-Pieri
products, conformal energies, and ground-state branching. It gives an
actual unitary model and corrects the entry's previously unresolved
status. We also describe a rank-fourteen unitary modular category with a
connected \'etale algebra whose module category is this rank-ten model.
An exact based-ring certificate specifies the identification with the
census tensor. These are identifications within established categorical
constructions; no claim of a new tensor category is made.
\end{abstract}
\maketitle

\section{The ring and its categorical model}

Let $R$ be the based ring with basis
\[
  1,g,g^2,g^3,g^4,x,gx,g^2x,g^3x,g^4x
\]
and relations
\begin{equation}\label{eq:ring}
 g^5=1,\qquad gx=xg,\qquad x^*=x,\qquad
 x^2=1+(g+g^2+g^3+g^4)x.
\end{equation}
The involution sends $g$ to $g^{-1}$. Put $d=2+\sqrt5$; then
\begin{equation}\label{eq:dim}
 d^2=1+4d,\qquad \FPdim(R)=5(1+d^2)=50+20\sqrt5.
\end{equation}
Appendix~\ref{sec:certificate} identifies $R$ exactly with
\texttt{R10\_1\_189}.

Write $\C=\C(\mathfrak{sl}_4,6)$ for the unitary WZW modular category.
Its simple objects are denoted $\wt{abc}$, where $(a,b,c)$ are Dynkin
labels, $a,b,c\geq0$ and $a+b+c\leq6$. The conformal embedding induced by
the symmetric-square representation of $\mathfrak{sl}_4$ is
\begin{equation}\label{eq:embedding}
 \SU(4)_6\subset\SU(10)_1.
\end{equation}
This is a realized conformal inclusion, with vacuum branching object
\begin{equation}\label{eq:B}
 B=\wt{000}\oplus\wt{060}\oplus\wt{202}\oplus\wt{222};
\end{equation}
see \cite[p.~29 of the arXiv version]{BE} and
\cite[\S\S1.4.3--1.4.5]{CS}. In particular, $B$ carries an actual
connected commutative separable algebra structure. Its category of local
modules is braided equivalent to $\C(\mathfrak{sl}_{10},1)$.
These assertions concern the conformal inclusion and its unitary
extension, rather than the existence of an algebra inferred from a
numerical modular invariant.

Let $\D=\C_B$ be the unitary fusion category of right $B$-modules, with
relative tensor product. Write $F:\C\to\D$ for free induction and $U$ for
forgetting the module structure. The functor $F$ is tensor and central
\cite[\S\S3.3--3.4, equation~(17)]{DMNO}. Denote the invertible simple
local modules by $\tau_j$, $j\in\mathbb Z/10$, with
$\tau_i\tau_j=\tau_{i+j}$. We suppress the symbol $\otimes_B$ for
products in $\D$.

\begin{theorem}\label{thm:model}
Put $\beta=F(\wt{010})$, $g=\tau_2$ and $x=\tau_5\beta$. The ten objects
\[
  g^j,\quad g^jx\qquad (0\leq j<5)
\]
are pairwise nonisomorphic simple objects of $\D$ and are closed under
tensor products and duality. Their full replete semisimple subcategory
$\K$ is a unitary fusion category satisfying $K_0(\K)\cong R$ as based
rings. Thus \texttt{R10\_1\_189} is unitarily categorifiable.
\end{theorem}

The proof below needs only a small amount of ground-state branching; it
does not require a completeness assertion about an affine branching
table. The additional construction in Section~\ref{sec:fourteen}
identifies $\K$ with the neutral component of $\C_B$ for the natural
$\mathbb Z/4$ grading.

\section{A minimal branching argument}

Let $v=\wt{010}$, the exterior square of the defining four-dimensional
representation. The exterior-Pieri rule gives
\begin{align}
 v^2&=\wt{000}+\wt{101}+\wt{020},\label{eq:v2}\\
 v^3&=3\wt{010}+2\wt{111}+\wt{030}+\wt{200}+\wt{002}.
 \label{eq:v3}
\end{align}
All dominant summands encountered have level at most three, so these
classical products are also the level-six fusion products.

\begin{lemma}\label{lem:branch}
For $j\in\mathbb Z/10$,
\begin{align}
 \dim\Hom_{\C}(v^2,U\tau_j)&=\delta_{j,0},\label{eq:hom2}\\
 \dim\Hom_{\C}(v^3,U\tau_j)&=
 \begin{cases}1,&j\in\{1,3,7,9\},\\0,&\text{otherwise.}\end{cases}
 \label{eq:hom3}
\end{align}
\end{lemma}

\begin{proof}
Conformality of \eqref{eq:embedding} means equality of the two Sugawara
conformal vectors. The lowest conformal energy of the level-one
$\mathfrak{sl}_{10}$ module indexed by $j$ is
\[
 h_j=\frac{j(10-j)}{20},\qquad 0\leq j\leq9.
\]
An affine $\mathfrak{sl}_4$ highest-weight module of highest weight
$\lambda$ occurring in it must satisfy
\begin{equation}\label{eq:energy}
 h_\lambda=h_j+n,\qquad n\in\mathbb Z_{\geq0}.
\end{equation}
If $n=0$, its multiplicity is the ordinary finite-dimensional
restriction multiplicity in the ground-state representation
$\Lambda^j\mathbb C^{10}$. With long roots of squared length two,
$h_\lambda=(\lambda,\lambda+2\rho)/20$. The needed values are
\[
\begin{array}{c|rrrrrrr}
\lambda&\wt{000}&\wt{101}&\wt{020}&\wt{010}&\wt{111}
 &\wt{200},\wt{002}&\wt{030}\\ \hline
h_\lambda&0&2/5&3/5&1/4&3/4&9/20&21/20.
\end{array}
\]
The residues of the $h_j$ are
$0,9/20,4/5,1/20,1/5,1/4$. Consequently $\wt{101}$ and
$\wt{020}$ cannot occur in any $U\tau_j$, while $\wt{000}$ occurs
only in $U\tau_0$, with multiplicity one. This proves \eqref{eq:hom2}.

For \eqref{eq:v3}, the residue of $h_{\wt{111}}$ excludes that weight.
The only possible index for $\wt{010}$ is $j=5$, but
$h_5=5/4>h_{\wt{010}}=1/4$, contradicting \eqref{eq:energy}. The weights
$\wt{200},\wt{002}$ can occur only for $j=1,9$, at ground energy.
The defining ten-dimensional representation restricts as
$\operatorname{Sym}^2 V$, $V=\mathbb C^4$, or its dual according to
orientation. Thus precisely one of these two weights occurs, once, at
each of $j=1,9$.

Finally, $\wt{030}$ can occur only for $j=3,7$, again at ground energy.
The ordinary Schur-functor identity
\begin{equation}\label{eq:plethysm}
 \Lambda^3(\operatorname{Sym}^2 V)
 =\mathbb S_{(4,1,1)}V\oplus\mathbb S_{(3,3)}V
\end{equation}
has highest weights $\wt{301}$ and $\wt{030}$. The latter is self-dual
and occurs once; dualizing gives the same multiplicity for $j=7$.
This proves \eqref{eq:hom3}. Identity \eqref{eq:plethysm} is also checked
as a full weight-character identity in the accompanying certificate.
\end{proof}

\begin{proof}[Proof of Theorem~\ref{thm:model}]
The object $v$ is self-dual, so $\beta$ is self-dual. Rigidity, the tensor
property of $F$, and its adjunction to $U$ give
\begin{align*}
 \dim\Hom_{\D}(\beta,\tau_j\beta)
 &=\dim\Hom_{\D}(\beta^2,\tau_j)
 =\dim\Hom_{\C}(v^2,U\tau_j),\\
 \dim\Hom_{\D}(\beta^2,\tau_j\beta)
 &=\dim\Hom_{\D}(\beta^3,\tau_j)
 =\dim\Hom_{\C}(v^3,U\tau_j).
\end{align*}
Lemma~\ref{lem:branch} proves that $\beta$ is simple and its ten
translates $\tau_j\beta$ are pairwise nonisomorphic. It also gives one
copy of each of $\tau_1\beta,\tau_3\beta,\tau_7\beta,\tau_9\beta$ in
$\beta^2$. Self-duality gives one copy of the unit.

For $[r]=\sin(r\pi/10)/\sin(\pi/10)$ the quantum Weyl formula gives
\[
 \FPdim(\beta)=\FPdim(v)=\frac{[4][3]}{[2]}=2+\sqrt5=d.
\]
The tensor functor $F$ preserves Frobenius--Perron dimension. The known
constituents of $\beta^2$ have total dimension $1+4d=d^2$; positivity
therefore proves the exact decomposition
\begin{equation}\label{eq:beta2}
 \beta^2=\one+\tau_1\beta+\tau_3\beta+\tau_7\beta+\tau_9\beta.
\end{equation}
Centrality of $F$ supplies isomorphisms $\beta\tau_j\cong\tau_j\beta$.
Hence $x=\tau_5\beta$ is self-dual, $g=\tau_2$ has order five, and
\eqref{eq:beta2} becomes \eqref{eq:ring}. The five invertibles and the
five indicated translates of $\beta$ are distinct, their dimensions
being $1$ and $d>1$. Relations \eqref{eq:ring} show closure under
products and duals. The full subcategory they span is consequently a
unitary fusion subcategory of $\D$, as asserted.
\end{proof}

\section{An accompanying rank-fourteen modular category}
\label{sec:fourteen}

The same conformal inclusion supplies a concrete modular category and
an algebra which recover $\K$. This strengthens the existence statement
beyond consistency of candidate modular data.

\begin{proposition}\label{prop:fourteen}
Let $J=\wt{600}$, let $L=\one\oplus J^2$, and put $\Q=\C_L^0$.
The image of $J$ in $\Q$ generates a semion subcategory $\mathcal S$.
Its M\"uger centralizer $\E=\mathcal S'\subset\Q$ is a unitary modular
category of rank fourteen and dimension $100d^2$. There is a connected
\'etale algebra
\[
 A=\one\oplus U_0\in\E,\qquad \FPdim(U_0)=d^2,
\]
obtained from the inclusion $L\subset B$, and a tensor equivalence
\[
 \E_A\simeq\K\simeq(\C_B)_0.
\]
Here the subscript $0$ denotes the neutral component of the
$\mathbb Z/4$ grading by $a+2b+3c$ modulo four.
\end{proposition}

\begin{proof}
The simple current $J$ has order four, with $J^2=\wt{060}$,
$\theta_J=i$ and $\theta_{J^2}=1$. We first justify the algebra
inclusion $L\subset B$. The regular right $B$-module is simple because
$B$ is connected. Tensoring it by the invertible $J^2$ again gives a
simple right module. Multiplication induces a nonzero right-module map
$J^2\otimes B\to B$ and hence an isomorphism. Its restriction to the
summand $J^2\otimes J^2=\one$ is nonzero and lands in the unique unit
summand of $B$. Thus multiplication restricts to $\one\oplus J^2$.
The restricted commutative algebra is the separable order-two group
algebra object; in particular $L$ is connected \'etale.

The image of $J$ in $\C_L^0$ is invertible of order two and twist $i$.
It generates a nondegenerate semion subcategory. Since $\Q$ is unitary
modular, its centralizer $\E$ is unitary modular as well.
To enumerate its simples, define the charge of $\wt{abc}$ by
$t=a+2b+3c\pmod4$. All four summands of $B$ have charge zero, so this
grading descends to $\C_B$. Simple-current monodromy with $J$ is
$\exp(\pm2\pi i t/4)$, where the sign depends on the braiding
convention. Locality for $L$ requires even charge; centralizing the
image of $J$ requires charge zero.

On weights of charge zero the action of $J^2$ is
\[
 (a,b,c)\longmapsto(c,6-a-b-c,a).
\]
There are twenty-two such weights. The only fixed weights are
$\wt{121}$ and $\wt{303}$. A nonfixed orbit gives one simple local
$L$-module. Each fixed weight gives two: its stabilizer is cyclic of
order two, and its twisted group algebra has two one-dimensional
irreducible representations. The resulting rank is
$(22-2)/2+2\cdot2=14$.

For completeness, Table~\ref{tab:fourteen} gives all dimensions and
twists. They follow from the quantum Weyl formula
\[
 \dim(\wt{abc})=
 \frac{[a+1][b+1][c+1][a+b+2][b+c+2][a+b+c+3]}{[2]^2[3]}
\]
and $\theta_\lambda=\exp(2\pi i h_\lambda)$. For a nonfixed orbit the
module dimension is the dimension of a representative; at a fixed
weight each of the two resolved modules has half that dimension.
Their squared dimensions sum to $100d^2$.

\begin{table}[ht]
\centering
\caption{Simple objects of $\E$, with $\zeta=\exp(2\pi i/5)$. A row
with two representatives describes two distinct orbits; a fixed-weight
row describes its two resolved modules.}
\label{tab:fourteen}
\begin{tabular}{@{}lccc@{}}
\toprule
Orbit representative(s)&Number&Dimension&Twist\\
\midrule
$\wt{000}$&1&1&1\\
$\wt{004},\wt{024}$&2&$1+d$&$\zeta$\\
$\wt{303}$ (fixed)&2&$1+d$&$\zeta^{-1}$\\
$\wt{012},\wt{032}$&2&$1+3d$&$\zeta^{-1}$\\
$\wt{121}$ (fixed)&2&$1+3d$&$\zeta$\\
$\wt{020}$&1&$2d$&$\zeta^{-2}$\\
$\wt{101}$&1&$2d$&$\zeta^2$\\
$\wt{105}$&1&$d$&1\\
$\wt{113}$&1&$5d$&$-1$\\
$\wt{202}$&1&$d^2$&1\\
\bottomrule
\end{tabular}
\end{table}

The algebra $B$ descends along $L\subset B$ to a connected \'etale
algebra $A$ in $\Q$ by \cite[Proposition~3.16]{DMNO}. Its underlying
object is $\one\oplus U_0$, where $U_0$ is the module of the orbit
$\{\wt{202},\wt{222}\}$. Both summands have charge zero, so $A$ lies
in $\E$ and has dimension $1+d^2$.

We detail the module-in-stages identification to fix its scope. Giving
an $A$-module in $\C_L^0$ amounts to giving a $B$-module whose restriction
to $L$ is local: the action maps are the same after balancing over $L$,
and the relative tensor products agree. A $B$-module supported on
charge-zero objects restricts to a local $L$-module and centralizes
the residual semion. Conversely, any $A$-module in $\E$ has underlying
$\C$-object of charge zero. Thus this tensor equivalence identifies
$\E_A$ with precisely $(\C_B)_0$.

Under the symmetric-square embedding the central element $iI_4$ maps
to $-I_{10}$. It therefore acts on the entire affine module $\tau_j$
by $(-1)^j$: its ground state has this character and the current modes
preserve it. Thus $U\tau_j$ has charge $2j\pmod4$.
Consequently $\tau_2$ has charge zero, whereas
$\tau_5$ and $\beta$ both have charge two. It follows that
$\K\subset(\C_B)_0\simeq\E_A$. Finally,
\[
 \FPdim(\E_A)=\frac{100d^2}{1+d^2}=5(1+d^2)=\FPdim(\K),
\]
using $(1+d^2)^2=20d^2$. The full fusion-subcategory inclusion is
therefore an equivalence.
\end{proof}

\begin{remark}
The proposition constructs both $\E$ and its algebra $A$ inside known
unitary categories. Table~\ref{tab:fourteen} records their dimensions
and twists; it is not used to assert uniqueness of a modular category
or to identify a full $S$-matrix from those invariants alone. The
rank-ten existence proof in Theorem~\ref{thm:model} is independent of
this optional modular construction.
\end{remark}

\appendix
\section{Exact identification with the census tensor}
\label{sec:certificate}

Let $(e_0,\ldots,e_9)$ be the original, zero-indexed census basis of
\texttt{R10\_1\_189}. The based-ring map sends the ordered basis in
Section~1 to
\begin{equation}\label{eq:permutation}
 (e_0,e_1,e_8,e_7,e_2,e_9,e_4,e_6,e_5,e_3).
\end{equation}
In particular, $g\mapsto e_1$ and $x\mapsto e_9$. If
$p=(0,1,8,7,2,9,4,6,5,3)$, the accompanying verifier checks
\[
 N^{\mathrm{census}}_{p(i),p(j)}{}^{p(k)}
 =N^{R}_{i,j}{}^k\qquad(0\leq i,j,k<10)
\]
for all $1000$ structure constants. The right-hand tensor is generated
directly from \eqref{eq:ring}. The source tensor is bound by SHA-256
of its compact JSON serialization, namely
\begin{center}
\small\ttfamily
c44e21d1844ded91dba30cd6c691180afa4d4f1be1c4f27806bbc12c8c8292f4.
\end{center}
This binds the certificate to the tensor itself rather than to a
potentially changing display label.

The standard-library Python files accompanying this note have the
following roles:
\begin{itemize}
\item \path{verify_conformal_model.py} checks the two Pieri products,
the needed intersections with the published branching table, dimension
exhaustion, and the exact census map \eqref{eq:permutation}; its receipt
is \path{conformal_model_certificate.json}.
\item \path{verify_minimal_branching.py} checks the conformal energies
over the rationals and verifies \eqref{eq:plethysm} as an equality of
four-variable weight characters by enumerating the exterior-cube and
semistandard-tableau characters. Its receipt is
\path{minimal_branching_certificate.json}.
\item \path{verify_r189_rank14_wzw.py} checks the charge-zero orbits,
fixed-point count, dimensions and twists in Table~\ref{tab:fourteen},
using exact arithmetic in
$\mathbb Q[t]/(t^4-5t^2+5)$ with $t=2\cos(\pi/10)$.
Its receipt is \path{r189_rank14_wzw_receipt.json}.
\end{itemize}
The finite checks supplement the categorical arguments. Existence of
the associativity constraints and of the algebra multiplications comes
from the realized conformal inclusion and the module constructions
above.

\begin{thebibliography}{9}

\bibitem{BE}
J.~B\"ocken-hauer and D.~E.~Evans,
\emph{Modular invariants from subfactors},
in \emph{Quantum Symmetries in Theoretical Physics and Mathematics},
Contemp. Math. \textbf{294}, Amer. Math. Soc., 2002, pp.~95--131.
\href{https://arxiv.org/abs/math/0006114}{arXiv:math/0006114v2}.

\bibitem{CS}
R.~Coquereaux and G.~Schieber,
\emph{Quantum symmetries for exceptional $\SU(4)$ modular invariants
associated with conformal embeddings},
SIGMA \textbf{5} (2009), 044, 31~pp.
\href{https://doi.org/10.3842/SIGMA.2009.044}{doi:10.3842/SIGMA.2009.044}.
\href{https://arxiv.org/abs/0805.4678}{arXiv:0805.4678}.

\bibitem{DMNO}
A.~Davydov, M.~M\"uger, D.~Nikshych and V.~Ostrik,
\emph{The Witt group of non-degenerate braided fusion categories},
J. Reine Angew. Math. \textbf{677} (2013), 135--177.
\href{https://doi.org/10.1515/crelle.2012.014}{doi:10.1515/crelle.2012.014}.
\href{https://arxiv.org/abs/1009.2117}{arXiv:1009.2117v2}.

\end{thebibliography}
\end{document}
